\section{Criterio unilateral de ventanas}
\label{sec:ventanas}

Esta nota completa el paso analítico que quedó abierto en la
Nota~\ref{sec:sumas}. Demostramos por qué una cota superior sobre intervalos
completos, disponibles tan lejos como se quiera, basta para excluir polos
de una transformada. La prueba conserva el signo de la función, controla
las colas de un núcleo y selecciona un polo mediante distancias en discos
compactos.

\fuente{A, Teorema~2.2, Lemas~2.3--2.4 y Corolario~2.5, pp.~4--6;
B, Proposición~2.1, Lema~2.2 y demostración, pp.~14--18. Al final se
documenta el cotejo con declaraciones Lean en una revisión fijada.
La demostración se reconstruye en lenguaje matemático; no se presenta
como una nueva ejecución de Lean.}

\subsection{Enunciado y orden de los cuantificadores}

Sea $g:[0,\infty)\to\R$ medible y supongamos que existe $C_0\ge0$ tal que
\begin{equation}
 |g(u)|\le C_0e^{2u}\qquad(u\ge0).
 \label{vent-crecimiento}
\end{equation}
Su transformada de Laplace, definida inicialmente para $\Ree s>2$, es
\[
 \Phi(s)=\int_0^\infty g(u)e^{-su}\dd u.
\]
Suponemos que admite continuación meromorfa a $\Ree s>0$, que todos sus
polos son simples y que el residuo de cada polo real es un número real
positivo. Denotamos por $Z$ su conjunto de polos. La meromorfía implica
que cada compacto contenido en ese semiplano contiene finitos polos.

Fijemos $0<\theta<1$ y $\delta>0$. La \emph{hipótesis integrada de ventanas}
significa lo siguiente: para cada $\eta>0$ y cada $U_0\ge1$, existe
$U\ge U_0$ y una función $R_U$ sobre $W_U=[U,(1+\delta)U]$ tales que
\[
 R_U(u)\ge0,\qquad g(u)\le R_U(u)\quad(u\in W_U),
\]
el producto $R_U(u)e^{-\theta u}$ es integrable en $W_U$ y
\begin{equation}
 \int_U^{(1+\delta)U}R_U(u)e^{-\theta u}\dd u\le e^{\eta U}.
 \label{vent-hipotesis}
\end{equation}
La ventana y su mayorante pueden depender de $\eta$ y del umbral.
Los parámetros $\theta,\delta$ y la constante $C_0$ permanecen fijos.
Exigir $U_0\ge1$ es equivalente a permitir cualquier umbral real: basta
reemplazarlo por $\max\{1,U_0\}$. Las integrales de esta nota son integrales
de Lebesgue; sus valores no cambian al añadir o quitar extremos de intervalos.

\begin{proposicion}[Criterio unilateral]
\label{vent-criterio}
Bajo estas hipótesis, todo polo $\rho\in Z$ satisface
$\Ree\rho\le\theta$.
\end{proposicion}

Esta formulación sigue la interfaz general \texttt{PoleData} junto con
\texttt{WindowHypU} de Lean. A añade la localización de los polos en
$\Ree s\le1$ y la regularidad en $s=1$, propias de su aplicación.
Nuestra prueba solo necesita que los polos queden a la izquierda de un
centro $\sigma>2$; esto se deduce de la propia representación de Laplace.
B parte directamente de una cota puntual, que obtendremos como corolario.

\subsection{Derivadas normalizadas y la desigualdad unilateral}

Para $a>0$ y todo entero $k\ge0$, integración por partes da
\begin{equation}
 \int_0^\infty\frac{u^ke^{-au}}{k!}\dd u=a^{-k-1}.
 \label{vent-integral-gamma}
\end{equation}
En efecto, la integral sin el factorial vale $1/a$ si $k=0$, y la
recurrencia es $I_k=(k/a)I_{k-1}$. Los términos de borde se anulan
porque la exponencial decreciente domina cualquier potencia.

Si $\Ree z_0>2$, ponemos $b_0=(\Ree z_0-2)/2>0$.
En el disco $|z-z_0|<b_0$, cada derivada del integrando de orden $j$
tiene módulo a lo sumo $C_0u^je^{-b_0u}$, integrable por
\eqref{vent-integral-gamma}. Esta dominación, también para el orden
siguiente, justifica diferenciar bajo el signo integral. Por tanto,
$\Phi$ es holomorfa en $\Ree s>2$ y
\begin{equation}
 A_k(z):=\frac{(-1)^k}{k!}\Phi^{(k)}(z)
 =\int_0^\infty g(u)\frac{u^ke^{-zu}}{k!}\dd u.
 \label{vent-derivadas}
\end{equation}
En particular, no hay polos con parte real mayor que $2$.
Todo polo queda estrictamente a la izquierda de cualquier $\sigma>2$.
Las identidades integrales se utilizan en este semiplano; más a la
izquierda se utiliza la continuación meromorfa.

Como $g$ es real, $A_k(\sigma)$ es real. Escribiendo
$g_+(u)=\max\{g(u),0\}$, la desigualdad triangular da
\begin{equation}
 \begin{aligned}
 |A_k(\sigma+it)|+A_k(\sigma)
 &\le\int_0^\infty (|g(u)|+g(u))
           \frac{u^ke^{-\sigma u}}{k!}\dd u\\
 &=2\int_0^\infty g_+(u)\frac{u^ke^{-\sigma u}}{k!}\dd u.
 \end{aligned}
 \label{vent-unilateral}
\end{equation}
La segunda contribución del lado izquierdo conserva su signo. Esta
identidad permite utilizar un mayorante de $g$ en la ventana sin pedir
allí una mejora de la cota para $|g|$. Fuera de la ventana sigue disponible
el crecimiento elemental \eqref{vent-crecimiento}.

{\raggedright
Véanse A, (2.9)--(2.10), p.~5; B, p.~17; y las declaraciones
\texttt{iteratedDeriv\_laplace} y \texttt{landau\_window\_ineq}.
\par}

\subsection{Concentrar el núcleo en una ventana completa}

\begin{lema}[Control uniforme de las colas]
\label{vent-colas}
Para cada $\delta>0$ existen $\sigma>2$ y $C_\delta>0$, independientes
de $U$, tales que, para todo $U\ge1$, se puede elegir un entero $k\ge U$
con
\begin{equation}
 \int_{[0,\infty)\setminus W_U}
 \frac{u^ke^{-(\sigma-2)u}}{k!}\dd u
 \le C_\delta\sigma^{-k}.
 \label{vent-cota-colas}
\end{equation}
Además, la elección satisface $k\asymp_\delta U$.
\end{lema}

\begin{proof}
Para $s>\ell>0$ y $0<U\le V$, separar una exponencial y aplicar
\eqref{vent-integral-gamma} proporciona
\begin{equation}
 \int_{[0,\infty)\setminus[U,V]}
 \frac{u^ke^{-su}}{k!}\dd u
 \le e^{-\ell V}(s-\ell)^{-k-1}
       +e^{\ell U}(s+\ell)^{-k-1}.
 \label{vent-colas-elementales}
\end{equation}
En la cola superior se usa $e^{-\ell u}\le e^{-\ell V}$; en la inferior,
$e^{\ell u}\le e^{\ell U}$. Las integrales restantes se amplían a toda
la semirrecta positiva.

Tomemos las constantes explícitas
\[
 \ell=2+\frac6\delta,\qquad b=1+\frac\delta2,\qquad
 M=\ell(1+\delta),\qquad \sigma=M(\ell+2),
\]
y el entero $k=\lceil b\sigma U\rceil$. Entonces $\sigma-2>\ell$ y
\begin{equation}
 b(\ell+2)=M-1,\qquad b(\ell-2)=\ell+1,\qquad
 b\sigma=M(M-1).
 \label{vent-identidades}
\end{equation}
La definición del techo implica
\begin{equation}
 \begin{gathered}
 b\sigma U\le k\le b\sigma U+1,\\
 U\le k\le(b\sigma+1)U\qquad(U\ge1).
 \end{gathered}
 \label{vent-orden-k}
\end{equation}

Pongamos $a_-=\sigma-2-\ell>0$ y $a_+=\sigma-2+\ell>0$.
La desigualdad $\log(\sigma/a)\le(\sigma-a)/a$ muestra que, para $a>0$,
\begin{equation}
 x+k\frac{\sigma-a}{a}\le c
 \quad\Longrightarrow\quad
 e^xa^{-k-1}\le\frac{e^c}{a\sigma^k}.
 \label{vent-comparacion-exponencial}
\end{equation}
Para el primer término de \eqref{vent-colas-elementales}, con
$s=\sigma-2$ y $V=(1+\delta)U$, las identidades
\eqref{vent-identidades} dan
\[
 -\ell(1+\delta)U+k\frac{\sigma-a_-}{a_-}
 =-MU+\frac{k}{M-1}\le\frac1{M-1}\le1.
\]
Para el segundo, observamos que $\ell a_+\le(\ell+1)\sigma$, pues esta
desigualdad equivale a $\ell(\ell-2)\le\sigma$, que se sigue de
$\sigma=\ell(1+\delta)(\ell+2)$. Así,
\[
 \ell U+k\frac{\sigma-a_+}{a_+}
 =\ell U-\frac{k(\ell-2)}{a_+}
 \le\ell U-\frac{(\ell+1)\sigma U}{a_+}\le0.
\]
Aplicar \eqref{vent-comparacion-exponencial} a ambas colas prueba
\eqref{vent-cota-colas}, con
\begin{equation}
 C_\delta=\frac e{\sigma-2-\ell}+\frac1{\sigma-2+\ell}.
 \label{vent-constante-colas}
\end{equation}
Todas las constantes se fijaron a partir de $\delta$, antes de conocer
la ventana. Esto prueba el lema.
\end{proof}

El núcleo normalizado
$(\sigma-2)^{k+1}u^ke^{-(\sigma-2)u}/k!$ tiene integral $1$.
Por \eqref{vent-cota-colas}, su masa exterior está acotada por
\[
 C_\delta(\sigma-2)
 \left(\frac{\sigma-2}{\sigma}\right)^k.
\]
La concentración es exponencial en $k$. Esa tasa, y no solo que la masa
exterior tienda a cero, permitirá comparar el núcleo con la contribución
de un polo.

{\raggedright
Véanse A, Lema~2.3, pp.~4--5, y \texttt{landau\_tails}.
B, Lema~2.2, p.~16, obtiene el mismo tipo de control mediante momentos
exponenciales de una densidad gamma.
\par}

\subsection{La cota superior para las derivadas}

Fijemos $\sigma,C_\delta$ como en el Lema~\ref{vent-colas}
y escribamos
\[
 E=\sigma-\theta,\qquad 0<E<\sigma.
\]
Para $v\ge0$, la serie exponencial da $v^k/k!\le e^v$.
Tomando $v=Eu$, resulta
\begin{equation}
 \frac{u^ke^{-\sigma u}}{k!}\le E^{-k}e^{-\theta u}.
 \label{vent-maximo-nucleo}
\end{equation}
En una ventana admisible, $g_+\le R_U$; fuera de ella,
$g_+\le C_0e^{2u}$. Aplicamos \eqref{vent-maximo-nucleo},
\eqref{vent-hipotesis} y el lema de colas a \eqref{vent-unilateral}.
Para el entero $k=\lceil b\sigma U\rceil$ obtenemos
\begin{equation}
 |A_k(\sigma+it)|+A_k(\sigma)
 \le 2e^{\eta U}E^{-k}+2C_0C_\delta\sigma^{-k}
 \qquad(t\in\R).
 \label{vent-cota-superior}
\end{equation}
La misma cota vale simultáneamente para todos los centros de la recta
$\Ree z=\sigma$. No hemos fijado todavía $\eta$: lo elegiremos después
de localizar el polo que contradiga esta estimación.

{\raggedright
Véanse A, (2.11), p.~5; \texttt{landau\_kernel\_max},
\texttt{landau\_window\_ineq} y \texttt{landau\_window\_key}.
\par}

\subsection{Elegir un polo mediante discos compactos}

Supongamos que existe $\rho_0\in Z$ con $\Ree\rho_0>\theta$.
No suponemos que haya un polo cuya parte real sea máxima.

El disco cerrado $\overline B(\sigma,E)$ queda dentro de $\Ree s>0$,
porque su extremo izquierdo tiene parte real $\theta>0$.
Su intersección $F$ con $Z$ es finita. Si $F$ es vacía, ponemos $D=E$.
Si no, ponemos
\[
 D=\min_{\rho\in F}|\sigma-\rho|.
\]
En ambos casos $0<D\le E$ y $B(\sigma,D)$ no contiene polos.
Consideramos dos alternativas.

\paragraph{Alternativa I.}
Existe un polo $\rho_1$ tal que $\sigma-\Ree\rho_1<D$.

\paragraph{Alternativa II.}
No existe tal polo. Esta posibilidad requiere $F\ne\varnothing$, pues,
si $D=E$ y $F$ fuera vacía, el polo $\rho_0$ satisfaría la alternativa I.
Un elemento que alcance la distancia mínima verifica
\[
 D=|\sigma-\rho|\ge\sigma-\Ree\rho\ge D.
\]
La igualdad fuerza $\rho=\beta:=\sigma-D\in\R$.
El punto más cercano es único. Además,
\[
 D\le\sigma-\Ree\rho_0<E,
\]
de modo que $\beta>\theta$. Como hay finitos polos en el disco cerrado,
existe $L_\beta$ con $D<L_\beta\le E$ para el cual $\beta$ es el único
polo de $B(\sigma,L_\beta)$.

Esto demuestra la dicotomía. En la segunda alternativa ya tenemos un
polo real aislado; falta precisar la elección de centro en la primera.

\subsubsection*{Un desplazamiento vertical elimina los empates}

En la alternativa I, sean $h=\sigma-\Ree\rho_1>0$ y
$t_0=\operatorname{Im}\rho_1$. Elegimos $\varepsilon>0$ tan pequeño que
\[
 h+\varepsilon<D,\qquad D+\varepsilon<\sigma.
\]
El conjunto $F_1=Z\cap\overline B(\sigma+it_0,D+\varepsilon)$ es finito.
Si $|t-t_0|<\varepsilon$ y $z=\sigma+it$, entonces $|z-\rho_1|<D$,
mientras que cualquier polo fuera de $F_1$ está a distancia mayor que
$D$ de $z$. Por tanto, el polo más cercano debe pertenecer a $F_1$.

Para dos polos distintos $a,b\in F_1$, la diferencia de cuadrados de sus
distancias a $\sigma+it$ es una función afín de $t$, con pendiente
$-2(\operatorname{Im}a-\operatorname{Im}b)$.
Si sus partes imaginarias son distintas, hay a lo sumo un valor de
empate. Si coinciden, el empate exigiría
\[
 (\sigma-\Ree a)^2=(\sigma-\Ree b)^2.
\]
Ambas distancias horizontales son positivas, por lo que tendríamos
$a=b$, contradicción. Podemos evitar los finitos valores de empate
dentro de $(t_0-\varepsilon,t_0+\varepsilon)$.

Fijamos uno de esos $t$. Existe entonces un único polo más cercano
$\rho$, con
\[
 0<d:=|\sigma+it-\rho|<D\le E.
\]
La finitud del conjunto relevante permite elegir un radio $L$ con
$d<L\le D$ tal que $B(\sigma+it,L)$ contenga solo ese polo.
Tomamos un radio interior $r$ con $d<r<L$. Así, $\Phi$ es holomorfa
cerca del disco cerrado $\overline B(\sigma,r)$, y, tras quitar la
parte principal en $\rho$, también lo será cerca de
$\overline B(\sigma+it,r)$.

En esta alternativa el polo elegido es necesariamente no real: si fuera
real, su distancia a $\sigma+it$ sería al menos su distancia a $\sigma$,
que es mayor o igual que $D$, en contradicción con $d<D$.
Para la estimación posterior bastará que su residuo sea no nulo.
En toda la selección se minimizó sobre conjuntos finitos dentro de compactos.

{\raggedright
Véanse A, Lema~2.4 y demostración del Teorema~2.2, pp.~5--6;
B, pp.~17--18; \texttt{landau\_finite}, \texttt{landau\_dichotomy}
y \texttt{landau\_nearest}.
\par}

\subsection{Cauchy hace visible el polo}

En la alternativa I, escribamos $z=\sigma+it$ y sea $c\ne0$
el residuo de $\Phi$ en $\rho$. La función
\[
 h_\rho(w)=\Phi(w)-\frac c{w-\rho}
\]
se prolonga holomórficamente en $\rho$ y en todo $B(z,L)$.
Aplicando Cauchy en la circunferencia de radio interior $r<L$, obtenemos
constantes $M_z,M_\sigma\ge0$ independientes de $k$ tales que
\begin{equation}
 \left|A_k(z)-\frac c{(z-\rho)^{k+1}}\right|\le M_zr^{-k},
 \qquad |A_k(\sigma)|\le M_\sigma r^{-k}.
 \label{vent-cauchy}
\end{equation}
Por ejemplo, $M_z$ es el máximo de $|h_\rho|$ en $|w-z|=r$;
ese máximo es finito por continuidad y compacidad.
El factor del polo procede de
\[
 \frac{\mathrm d^k}{\mathrm dw^k}\frac c{w-\rho}
 =\frac{(-1)^kk!c}{(w-\rho)^{k+1}}.
\]
La desigualdad triangular inversa y la realidad de $A_k(\sigma)$ dan
\begin{equation}
 |A_k(z)|+A_k(\sigma)
 \ge\frac{|c|}{d^{k+1}}-(M_z+M_\sigma)r^{-k}.
 \label{vent-cota-polo}
\end{equation}
El valor absoluto elimina la fase compleja de la contribución dominante;
basta que $c\ne0$.

En la alternativa II ponemos $d=\sigma-\beta$ y elegimos
$d<r<L_\beta$. El residuo $c_\beta$ es real positivo.
Al quitar la parte principal y aplicar Cauchy se obtiene
\begin{equation}
 A_k(\sigma)\ge\frac{c_\beta}{d^{k+1}}-Mr^{-k}.
 \label{vent-polo-real}
\end{equation}
La expresión de la izquierda de \eqref{vent-cota-superior}, ahora
con $t=0$, es al menos $A_k(\sigma)$. Por tanto, en ambas alternativas
existen constantes $a>0$, $M\ge0$ y distancias $0<d<r<E$ tales que
\begin{equation}
 |A_k(\sigma+it)|+A_k(\sigma)
 \ge a\,d^{-k}-Mr^{-k}
 \qquad(k\ge0).
 \label{vent-cota-inferior}
\end{equation}
En la primera alternativa $a=|c|/d$; en la segunda $a=c_\beta/d$.
Los centros, radios y constantes ya están fijados antes de elegir
una ventana. El radio interior garantiza que Cauchy se aplica sobre
una circunferencia contenida en la región de holomorfía.

{\raggedright
Véanse A, p.~6; B, pp.~17--18; \texttt{landau\_cauchy}
y \texttt{landau\_pole\_asym}.
\par}

\subsection{Contradicción y papel del residuo positivo}

\begin{proof}[Conclusión de la demostración de la Proposición~\ref{vent-criterio}]
Fijemos ahora
\[
 \eta=\frac12\log(E/d)>0.
\]
La hipótesis de ventanas proporciona, para este $\eta$, valores de $U$
arbitrariamente grandes y los correspondientes $k\ge U$.
Al combinar \eqref{vent-cota-superior} y \eqref{vent-cota-inferior}
y multiplicar por $d^k$, resulta
\begin{equation}
 a\le M\left(\frac dr\right)^k
 +2e^{\eta U}\left(\frac dE\right)^k
 +2C_0C_\delta\left(\frac d\sigma\right)^k.
 \label{vent-contradiccion}
\end{equation}
Como $U\le k$,
\[
 e^{\eta U}\left(\frac dE\right)^k
 \le\left(e^\eta\frac dE\right)^k
 =\left(\frac dE\right)^{k/2}.
\]
Las tres bases exponenciales son estrictamente menores que $1$.
Al tomar ventanas con $U\to\infty$, también $k\to\infty$, y el lado
derecho de \eqref{vent-contradiccion} tiende a cero.
Esto contradice $a>0$ y demuestra la proposición.
\end{proof}

El orden de elección fue: datos $g,\theta,\delta$; constantes del núcleo;
un polo hipotético, centros, radios y constantes de Cauchy; el margen
$\eta$; finalmente una ventana suficientemente lejana y su orden $k$.
No hace falta controlar los huecos entre ventanas ni disponer de la
estimación para todos los órdenes de derivación.
Véanse A, p.~6; \texttt{landau\_window} y
\texttt{landau\_growth\_contra}.

La hipótesis de signo sobre los polos reales es necesaria.
Si $0<\theta<\beta<1$ y $g(u)=-e^{\beta u}$, entonces $g(u)\le0$
en toda ventana, pero
\[
 \Phi(s)=-\frac1{s-\beta}
\]
tiene un polo en $\beta>\theta$ con residuo $-1$.
En ese ejemplo, $A_k(\sigma)<0$ y
$|A_k(\sigma)|+A_k(\sigma)=0$. El residuo positivo de la
alternativa II impide precisamente esta cancelación.

\subsection{La versión puntual y su aplicación a las funciones L}

\begin{proposicion}[Corolario: cota puntual]
\label{vent-corolario}
Supongamos las hipótesis analíticas de la Proposición~\ref{vent-criterio}.
Si hay una constante fija $C\in\R$ y valores arbitrariamente grandes
de $Y$ tales que
\begin{equation}
 H(x)\le Cx^\theta
 \qquad\text{para todo }x\in[Y,Y^{1+\delta}],
 \label{vent-cota-puntual}
\end{equation}
donde $g(u)=H(e^u)$, entonces $\Phi$ no tiene polos con parte real
mayor que $\theta$.
\end{proposicion}

\begin{proof}
Reemplazamos $C$ por $C_+=\max\{C,0\}$ y ponemos $U=\log Y$.
La ventana se transforma en $[U,(1+\delta)U]$ y
$g(u)\le C_+e^{\theta u}$. Tomando $R_U(u)=C_+e^{\theta u}$,
\[
 \int_U^{(1+\delta)U}R_U(u)e^{-\theta u}\dd u=C_+\delta U.
\]
Para cada $\eta>0$, esta cantidad es a lo sumo $e^{\eta U}$ cuando
$U$ es suficientemente grande. Por ejemplo, basta
$U\ge\max\{1,2C_+\delta/\eta^2\}$, pues
$e^{\eta U}\ge\eta^2U^2/2$.
Como las ventanas son arbitrariamente lejanas, se puede imponer
ese umbral. Aplicamos la Proposición~\ref{vent-criterio}.
\end{proof}

Hay también una ruta directa, utilizada por B: dentro de la ventana,
\eqref{vent-integral-gamma} proporciona
\begin{equation}
 \begin{aligned}
 |A_k(\sigma+it)|+A_k(\sigma)
 &\le 2C_+E^{-k-1}+2C_0C_\delta\sigma^{-k}\\
 &\le\left(\frac{2C_+}{E}+2C_0C_\delta\right)E^{-k}.
 \end{aligned}
 \label{vent-ruta-puntual}
\end{equation}
La comparación con \eqref{vent-cota-inferior} conduce a la misma
contradicción sin introducir $\eta$. La versión integrada permite
cotejar el enunciado formal más general.
Véanse A, Corolario~2.5, p.~6, y B, Proposición~2.1 y (2.4),
pp.~14--18.

Para una familia finita de caracteres primitivos y pesos $a_\chi>0$,
la construcción de la Nota~\ref{sec:sumas} da
\[
 H(x)=\sum_\chi a_\chi\Ree G_\chi(x),\qquad
 \Phi(s)=\frac12\sum_\chi a_\chi
 \bigl(F_\chi(s)+F_{\bar\chi}(s)\bigr).
\]
El crecimiento $|g(u)|\le C_0e^{2u}$ procede de la cota elemental de las
sumas suavizadas. La continuación meromorfa y la regularidad en $1$
proceden de sus fórmulas de Mellin. En un cero representado por la
familia, el residuo es
\begin{equation}
 \begin{gathered}
 \Res_{s=\rho}\Phi(s)=\frac{b_\rho}{\rho(1-\rho)},\\
 b_\rho=\frac12\sum_\chi a_\chi
 \bigl(m_\chi(\rho)+m_{\bar\chi}(\rho)\bigr)>0.
 \end{gathered}
 \label{vent-residuo-familia}
\end{equation}
Los pesos positivos evitan la cancelación de polos. En un polo real
$\beta\in(0,1)$, también dan el signo positivo del residuo.
Para situar los polos reales en ese intervalo bastan el producto
de Euler y la regularidad en $1$; no se supone previamente la
hipótesis de Riemann.

Aplicamos el corolario a $H_\psi$ y $\delta=1/2$.
Si existe $\varepsilon_0>0$ tal que la construcción aritmética entrega
\eqref{vent-cota-puntual} con $\theta=1/2+\varepsilon$ para cada
$0<\varepsilon<\varepsilon_0$, cualquier cero con parte real
$\beta>1/2$ queda excluido al elegir
\[
 0<\varepsilon<\min\{\varepsilon_0,\,\beta-1/2,\,1/2\}.
\]
La ecuación funcional permite tratar después el lado izquierdo de la
franja. Para caracteres no primitivos se retiran los factores de Euler
adicionales, no nulos en $\Ree s>0$. Estas últimas reducciones se
identifican en A, pp.~6--7, y B, Proposición~2.3, pp.~18--19.
La construcción aritmética de las ventanas sigue pendiente de desarrollo
en el cuaderno: haber demostrado el criterio no suministra esa estimación.

\subsection{Concordancia con Lean y alcance del cotejo}

Se consultó el espejo
\href{https://github.com/CocoSgt/riemann-hypothesis-lean/tree/7cbf7a0c32fca1badf05931879a37c3f8a2b03b4}
{\texttt{CocoSgt/riemann-hypothesis-lean}},
en la revisión
\href{https://github.com/CocoSgt/riemann-hypothesis-lean/tree/7cbf7a0c32fca1badf05931879a37c3f8a2b03b4}
{\texttt{7cbf7a0c32fc}},
que declara Lean \texttt{4.33.1} y Mathlib
\texttt{0df444a360ea}. Los identificadores completos se conservan en
el registro del repositorio. Esta referencia fija el código leído;
no afirma que hayamos comparado todo el espejo con el archivo original
distribuido por Alpöge.

\begingroup
\raggedright
\begin{description}
 \item[Datos analíticos.]
 \texttt{PoleData} incluye medibilidad, crecimiento, identidad de Laplace,
 holomorfía fuera de $Z$, ausencia de acumulación local y una parte
 principal simple con residuo real positivo cuando el polo es real.

 \item[Hipótesis integrada.]
 \texttt{WindowHypU} exige la integrabilidad del mayorante ponderado de
 \eqref{vent-hipotesis} y ventanas arbitrariamente lejanas para cada
 $\eta>0$.

 \item[Derivadas y núcleo.]
 \texttt{iteratedDeriv\_laplace} y \texttt{kernel\_integral} identifican
 las derivadas y la integral del núcleo. La representación integral
 de las derivadas se utiliza solo en $\Ree z>2$.

 \item[Control uniforme.]
 \texttt{landau\_tails}, \texttt{landau\_kernel\_max} y
 \texttt{landau\_window\_ineq} fijan las constantes antes de la ventana
 y conservan el signo del término real.

 \item[Geometría.]
 \texttt{landau\_pole\_re\_le}, \texttt{landau\_finite},
 \texttt{landau\_dichotomy} y \texttt{landau\_nearest} obtienen
 $\Ree\rho\le2$ y efectúan la selección en conjuntos finitos.

 \item[Estimaciones locales.]
 \texttt{landau\_cauchy} y \texttt{landau\_pole\_asym} restan la parte
 principal y aplican Cauchy con un radio estrictamente interior.

 \item[Conclusión.]
 \texttt{landau\_window\_key}, \texttt{landau\_window\_contra} y
 \texttt{landau\_window} comparan la cota de derivadas con la contribución
 de un polo de parte real mayor que $\theta$.

 \item[Puente desde cotas puntuales.]
 \texttt{windowHyp\_of\_rpow}, \texttt{windowHypU\_of\_windowHyp}
 y \texttt{zeroFree\_of\_rpow\_window}
 deducen la condición integrada de la cota puntual y la transportan
 a \texttt{WindowHypU} mediante el cambio $x=e^u$; el corolario
 concluye la no anulación.
\end{description}
\par
\endgroup

{\raggedright
Las definiciones están en
\href{https://github.com/CocoSgt/riemann-hypothesis-lean/blob/7cbf7a0c32fca1badf05931879a37c3f8a2b03b4/src/Development/AutomorphicArithmetic/DedekindLFunctions/LandauZeroFreeCriterion.lean}
{\texttt{LandauZeroFreeCriterion.lean}};
el ensamble final está en
\href{https://github.com/CocoSgt/riemann-hypothesis-lean/blob/7cbf7a0c32fca1badf05931879a37c3f8a2b03b4/src/Development/AutomorphicArithmetic/DedekindLFunctions/LandauTheorems/landauWindow.lean}
{\texttt{landauWindow.lean}}.
La definición llamada
\texttt{Deduction.}\allowbreak\texttt{LandauZeroFreeCriterion},
al final del primer archivo, expresa otra variante con una cota eventual
sobre una semirrecta; el teorema de ventanas estudiado aquí es
\texttt{Deduction.landau\_window}.
El registro del repositorio conserva las rutas, declaraciones y hashes
de los archivos recuperados.
\par}

El criterio general no incluye como hipótesis la ausencia de polos en
$\Ree s>1$ ni la regularidad en $1$; estas se satisfacen en la aplicación
a las funciones $L$. \texttt{WindowHypU} exige un mayorante puntual
no negativo de $g$ y controla su tamaño mediante la integral ponderada;
no exige una cota puntual de la forma $g(u)\le Ce^{\theta u}$.
A, B y Lean concuerdan en el mecanismo revisado, con esas diferencias
de formulación.

\paragraph{Estado del tramo.}
Se ha reconstruido la demostración matemática y se ha efectuado un
cotejo localizado con los enunciados y pruebas señalados.
Este cotejo no incluyó compilación ni ejecución de Lean, ni una auditoría
de todas sus dependencias. Esta nota cierra el pendiente analítico del
cuaderno; el siguiente desarrollo es la identidad de caracteres y
la cancelación de su término principal.
